% !TeX program = xelatex
\documentclass[11pt,a4paper]{ctexart}

% ========== Packages ==========
\usepackage[top=2cm,bottom=2cm,left=2cm,right=2cm]{geometry}
\usepackage{setspace}
\singlespacing
\usepackage[compact]{titlesec}
\titlespacing{\section}{0pt}{8pt}{4pt}
\titlespacing{\subsection}{0pt}{6pt}{2pt}
\titlespacing{\subsubsection}{0pt}{4pt}{2pt}
\usepackage{fancyhdr}
\usepackage{graphicx}
\usepackage{float}
\usepackage{amsmath,amssymb}
\usepackage{textcomp}
\usepackage{algorithm}
\usepackage{algpseudocode}
\usepackage{booktabs}
\usepackage{tabularx}
\setlength{\floatsep}{6pt plus 2pt minus 2pt}
\setlength{\textfloatsep}{8pt plus 2pt minus 2pt}
\setlength{\intextsep}{6pt plus 2pt minus 2pt}
\setlength{\abovecaptionskip}{4pt}
\setlength{\belowcaptionskip}{4pt}

\usepackage{hyperref}
\usepackage{enumitem}
\setlist{nosep,leftmargin=1.5em}
\usepackage{caption}
\captionsetup{font=small,skip=4pt}
\usepackage{subcaption}
\usepackage{tikz}
\usetikzlibrary{arrows.meta,shapes,positioning,fit,backgrounds}
\usepackage{listings}
\usepackage{xcolor}

\setlength{\parskip}{2pt plus 1pt minus 1pt}
\setlength{\parindent}{0pt}

% gbt7714 not needed — using manual thebibliography

% ========== Code listing style ==========
\lstdefinestyle{pystyle}{
    language=Python,
    basicstyle=\ttfamily\small,
    keywordstyle=\color{blue},
    stringstyle=\color{red},
    commentstyle=\color{gray},
    numbers=left,
    numberstyle=\tiny\color{gray},
    frame=single,
    breaklines=true,
    showstringspaces=false,
    captionpos=b
}

\lstdefinestyle{jsonstyle}{
    language=,
    basicstyle=\ttfamily\small,
    stringstyle=\color{red},
    frame=single,
    breaklines=true,
    showstringspaces=false
}

% ========== Header/Footer ==========
\pagestyle{fancy}
\fancyhf{}
\fancyhead[L]{\footnotesize Micro-MDT: Compute-Optimal 慢病复诊智能副驾}
\fancyhead[R]{\footnotesize \thepage}
\renewcommand{\headrulewidth}{0.3pt}
\setlength{\headheight}{14pt}
\setlength{\headsep}{8pt}

% ========== Title info ==========
\title{{\Huge \textbf{Micro-MDT}}\\[6pt]
       {\Large 基于 Compute-Optimal 多智能体辩论的}\\[4pt]
       {\Large 慢病复诊智能副驾系统}\\[12pt]
       {\normalsize Micro-MDT: A Compute-Optimal Multi-Agent Debate System}\\[4pt]
       {\normalsize for Chronic Disease Follow-up Copilot}}

\author{姓名A \ 学号XXXXXX \quad 姓名B \ 学号XXXXXX \quad 姓名C \ 学号XXXXXX \\[4pt]
        \small 指导教师：XXX 教授}
\date{\today}

\begin{document}

\maketitle
\thispagestyle{empty}

% ===================================================================
\newpage
\setcounter{page}{1}
% ===================================================================

% ==================== 中文摘要 ====================
\begin{center}
    {\large \textbf{摘\quad 要}}
\end{center}

\noindent {\bf 背景：} 大语言模型（LLM）在医疗领域的应用面临可靠性与责任边界不足的挑战。单一模型直接生成诊疗建议时，可能出现事实幻觉、药物相互作用遗漏和急症识别延迟等安全风险。Scaling Test-Time Compute 理论表明，针对复杂任务在推理阶段增加计算资源，可能比单纯扩大模型参数更具成本效益，这为面向医疗安全的多智能体工作流提供了理论动机。然而，如何在可控成本下将 Test-Time Compute、Multi-Agent Debate 与 Human-in-the-Loop 机制工程化为可复现实验系统，仍有待进一步探索。

\noindent {\bf 方法：} 本文提出并实现 Micro-MDT：一个基于 Compute-Optimal 策略的低成本多智能体医疗安全工作流原型。系统包含七个专门化 Agent：Triage Agent 负责病例分诊与算力路由；Generalist Agent 作为方案生成者（Proposer）；Pharmacist Agent 与 SafetyEthics Agent 作为双层验证器（Verifier）；GeneralistRevision Agent 执行顺序修正；DecisionSynthesis Agent 生成结构化人类判决选项；Documentation Agent 根据人类决策自动生成电子病历草稿、用药说明和复诊提醒。系统采用纯 Prompt Engineering 架构，无需模型微调，支持 Mock 确定性规则引擎和 OpenAI-Compatible API 两种 Provider 模式，并内置基于 stdlib HTTP Server 的零依赖交互式 Web 界面。

\noindent {\bf 实验：} 在包含 12 个手工构造病例的测试集上，本文覆盖 LOW（稳定复诊）、MEDIUM（轻度异常）和 HIGH（多重用药冲突/急症风险）三个难度等级，并对比单模型直接生成模式与 Micro-MDT 多智能体辩论模式。实验结果表明：(1) Mock 模式下分诊路由准确率为 83.3\%，误差主要来自关键词规则对稳定慢病语境的假阳性；(2) 6 个 HIGH 难度病例均触发安全拒答，拒答触发率为 100\%；(3) Human-in-the-Loop 回路能够在分歧场景下生成结构化 EHR 三件套文书草稿。所有病例均为模拟或手工构造文本，系统输出不构成真实医疗建议。

\noindent {\bf 结论：} Micro-MDT 表明，Compute-Optimal 多智能体辩论架构可作为高风险医疗 AI 安全机制的可复现实验平台。系统开箱即用、零额外依赖、兼容低成本 OpenAI-compatible API，为研究医疗 AI 的安全拒答、人机协同与推理期算力分配提供了参考实现。

\vspace{1em}
\noindent {\bf 关键词：} 多智能体系统；Test-Time Compute；医疗AI安全；安全拒答；人在回路；慢病复诊

\vspace{1em}
\noindent {\bf 中图分类号：} TP391 \qquad {\bf 文献标识码：} A

% ==================== 英文摘要 ====================
\newpage
\begin{center}
    {\large \textbf{Abstract}}
\end{center}

\noindent {\bf Background:} Large Language Models (LLMs) face reliability and accountability challenges in high-stakes medical applications. Single-model diagnostic outputs may suffer from hallucinations, missed drug interactions, and delayed emergency recognition. Compute-Optimal Test-Time Compute~\cite{snell2024scaling} suggests that allocating additional inference-time computation to difficult cases can be more cost-effective than uniformly scaling model parameters, motivating multi-agent architectures for medical AI safety. However, the practical integration of Proposer-Verifier debate, sequential revision, abstention, and Human-in-the-Loop control remains underexplored.

\noindent {\bf Methods:} We introduce Micro-MDT, a seven-agent medical safety workflow prototype that implements Compute-Optimal routing. Agents include Triage (difficulty-based routing), Generalist (Proposer), Pharmacist and SafetyEthics (dual Verifiers), Revision (sequential correction), DecisionSynthesis (structured human decision options), and Documentation (EHR-style document drafting). The system uses pure prompt engineering, requires no model fine-tuning, supports deterministic Mock and OpenAI-compatible API providers, and includes a zero-dependency interactive Web UI.

\noindent {\bf Results:} On a 12-case manually constructed test suite spanning LOW (stable follow-up), MEDIUM (mild abnormality), and HIGH (polypharmacy/emergency) tiers, Micro-MDT achieves 83.3\% triage routing accuracy in Mock mode, with errors attributable to rule-based false positives for stable chronic-disease cases. All six HIGH-risk cases trigger safety abstention, and the Human-in-the-Loop path successfully generates EHR-style document drafts. In a single-model baseline comparison, 66.7\% of critical safety risks are insufficiently flagged, whereas Micro-MDT intercepts all HIGH-risk cases through dual verification.

\noindent {\bf Conclusions:} Micro-MDT demonstrates the feasibility of a reproducible, zero-dependency prototype for studying Compute-Optimal multi-agent debate in medical AI safety. With weighted routing, the estimated per-case API cost is approximately \textyen 0.0027. The system is intended for research and educational evaluation only, not for real clinical decision-making.

\vspace{1em}
\noindent {\bf Keywords:} Multi-Agent Systems; Test-Time Compute; Medical AI Safety; Abstention; Human-in-the-Loop; Chronic Disease Follow-up

% ===================================================================
\newpage
% ==================== 第1章 引言 ====================
\section{引言}

\subsection{研究背景与问题}

慢性非传染性疾病（如糖尿病、高血压）是全球公共卫生的重要负担。据世界卫生组织统计，全球约有 4.22 亿糖尿病患者和 12.8 亿高血压患者 \cite{world2024diabetes}。慢病管理依赖持续复诊、用药评估与风险分层；然而，医疗资源分布不均使大量患者难以及时获得规范化随访服务。据《中国卫生健康统计年鉴》，2023 年全国每千人口执业医师数为 3.4 人 \cite{china2023health}，基层慢病管理中的医患供需矛盾仍较突出。

大语言模型（Large Language Models, LLMs）的快速发展为缓解这一矛盾提供了新的技术可能。GPT-4、DeepSeek、Qwen 等通用或医学增强模型在医学知识问答中展现出较强能力，部分基准测试得分已接近或超过人类专家水平。然而，医疗是典型的高风险决策领域：一次错误的用药建议可能直接危及患者生命。单一模型直接输出诊疗方案至少存在三个安全隐患：

\begin{enumerate}[label=(\arabic*),leftmargin=2em]
    \item \textbf{幻觉与事实错误}：LLM 可能生成看似合理但实际错误的诊断或用药建议，且表述的自信度与准确性之间无必然关联；
    \item \textbf{药物冲突遗漏}：在多重用药场景中，单次推理可能忽略药物相互作用、肝肾禁忌和过敏史等关键风险因素；
    \item \textbf{急症延误}：模型可能将需要紧急处理的危险信号（如胸痛、黑便）错误归类为可观察的轻症，造成致命延误。
\end{enumerate}

这些问题的根源在于，单次 LLM 推理通常缺乏外部化、可追踪的验证机制。人类医疗体系常通过多学科会诊（Multi-Disciplinary Team, MDT）降低复杂决策风险：不同专科医生从各自视角审查方案，通过交叉验证和讨论形成相对稳健的共识。这一模式为医疗 AI 工作流设计提供了重要启示。

\subsection{理论动机：Scaling Test-Time Compute}

Snell 等人于 2024 年提出的 Scaling Test-Time Compute 理论 \cite{snell2024scaling} 为多智能体医疗系统提供了重要方法论启发。该研究表明，在推理阶段增加计算资源（如并行采样、自我修正、搜索验证）在部分任务上可以比单纯增大模型参数规模更有效地提升性能，即 ``Compute-Optimal'' 策略：简单问题消耗少量算力直接输出，复杂问题消耗更多算力进行深度推理和验证。

将这一理论映射到多智能体医疗系统，自然形成了四个设计原则：

\begin{enumerate}[label=(\arabic*),leftmargin=2em]
    \item \textbf{Proposer vs. Verifier}：诊断 Agent 生成方案（Proposer），药学 Agent 和安全 Agent 充当验证器（Verifier），形成双层审查；
    \item \textbf{Sequential Revisions}：审查不通过时，将反馈交回 Proposer 进行修订，形成多轮辩论；
    \item \textbf{安全拒答（Abstention）}：多 Agent 的分歧作为触发拒答的显式信号，取代单模型的概率阈值判断；
    \item \textbf{算力最优分配（Compute-Optimal）}：通过分诊 Agent 评估病例难度，按需分配计算资源。
\end{enumerate}

\subsection{项目目标与贡献}

本文提出 Micro-MDT，一个基于 Compute-Optimal 策略的低成本多智能体慢病复诊智能副驾原型系统。本文的核心贡献如下：

\begin{itemize}[leftmargin=2em]
    \item[\textbf{C1.}] \textbf{七 Agent Compute-Optimal 工作流。} 将 Test-Time Compute 的 Proposer-Verifier、Sequential Revision 和 Abstention 机制工程化为可运行的 Python 系统，包含 Triage（路由）、Generalist（生成）、Pharmacist + SafetyEthics（双层审查）、Revision（修正）、DecisionSynthesis（决策选项）和 Documentation（文书生成）七个角色。
    \item[\textbf{C2.}] \textbf{三级算力自适应路由。} 低难度病例（$\sim$60\%）走 2 次 API 调用的快速通道；中等难度（$\sim$30\%）触发 1--2 轮交叉复核；高难度（$\sim$10\%）进入最多 3 轮 MDT 辩论，未收敛则触发安全拒答。在给定病例比例和 API 单价假设下，加权平均成本约 \textyen 0.0027/次。
    \item[\textbf{C3.}] \textbf{结构化 Human-in-the-Loop 回路。} AI 内部分歧时不直接输出诊疗结论，而是生成 A/B/C 结构化决策选项；人类医生给出判决后，系统自动触发 EHR 三件套文书草稿（复诊记录、用药指导、复诊提醒）生成。
    \item[\textbf{C4.}] \textbf{零依赖、开箱即用的完整系统。} 纯 Python stdlib 实现，无 LangChain/AutoGen 等第三方依赖；支持 Mock 确定性模式和 OpenAI-Compatible API 两种 Provider；内置基于 stdlib HTTP Server 的交互式 Web 界面。
\end{itemize}

\subsection{论文组织}

本文其余部分组织如下：第 2 章综述相关研究；第 3 章详述系统设计与方法；第 4 章展示实验结果与分析；第 5 章进行讨论；第 6 章给出结论与展望。

% ===================================================================
\newpage
% ==================== 第2章 文献综述 ====================
\section{文献综述}

\subsection{大语言模型在医疗领域的应用}

近年来，大语言模型在医疗领域取得了显著进展。Singhal 等人提出的 Med-PaLM 2 \cite{singhal2023medpalm} 在美国医学执照考试（USMLE）中达到了 86.5\% 的专家级准确率。OpenAI 的 GPT-4 在多种医学基准测试中展现出与专业医生相当的诊断能力 \cite{nori2023capabilities}。然而，这些评估主要集中于知识问答层面，距离可监管、可追责的临床安全部署仍有显著差距。

在安全性和可靠性方面，既有研究揭示了 LLM 在医疗场景中的系统性缺陷。Lee 等人 \cite{lee2023benefits} 指出，GPT-4 等通用模型虽可辅助医学信息处理，但在真实医疗环境中仍存在事实错误、责任归属和安全边界等风险。对于慢病复诊场景，药物相互作用、肾功能约束和急症信号识别尤其需要外部验证机制。

在慢病管理领域，LLM 可承担病史整理、风险摘要和患者教育等辅助任务，但直接生成用药调整建议仍需谨慎。本研究因此将系统定位为``慢病复诊智能副驾''，强调人机协同和安全拒答，而非端到端自动诊疗。

\subsection{多智能体辩论系统}

多智能体辩论（Multi-Agent Debate）是近期 AI 安全研究的热点方向。Du 等人 \cite{du2024improving} 提出了多智能体辩论框架，通过多个 LLM 实例之间的辩论来提升推理的事实性和逻辑性。实验表明，3-5 个 Agent 的多轮辩论可将数学推理准确率提升 15-25 个百分点。

Zhuoyun Du 等人 \cite{liang2024multiagent} 将多团队协作框架应用于软件开发场景，通过跨团队交流探索多个候选决策路径。Chan 等人 \cite{chan2024chateval} 提出的 ChatEval 框架利用多智能体辩论替代传统 NLP 评估指标，在文本质量评估中展现出与人类评判更高的一致性。

在医疗领域，Tang 等人 \cite{tang2024medagents} 初步探索了多 Agent 医疗诊断系统，设计了分诊、诊断和验证三个角色，在 MedQA 数据集上展示了优于单模型的性能。但其系统未引入安全拒答机制和 Human-in-the-Loop 回路。

\subsection{Test-Time Compute Scaling}

Snell 等人 2024 年的工作 \cite{snell2024scaling} 系统性地研究了 Test-Time Compute 的规模化规律。论文发现：(1) 对于简单问题，增加推理计算资源收益递减；(2) 对于复杂问题，通过并行采样和验证器搜索可显著提升性能；(3) Compute-Optimal 策略可在固定计算预算下最大化准确率。

这一方向也与 Chain-of-Thought、Self-Consistency 和 Tree of Thoughts 等推理增强方法相关。Wei 等人 \cite{wei2022chain} 表明，显式生成中间推理步骤可以提升复杂推理任务表现；Wang 等人 \cite{wang2023self} 进一步通过多路径采样与一致性投票提升答案稳健性；Yao 等人 \cite{yao2024tree} 则将推理过程组织为搜索树，为复杂问题提供了更结构化的探索机制。这些工作共同说明，推理期计算资源的组织方式会显著影响模型输出质量。

Lightman 等人 \cite{lightman2024let} 的研究表明，过程监督（Process Supervision）优于结果监督（Outcome Supervision）。这一发现支持了在复杂任务中显式审查中间推理过程的必要性，也为多 Agent 交叉审查比单模型自我检查更稳健的论点提供了方法论支撑。

\subsection{Human-in-the-Loop 医疗 AI}

Human-in-the-Loop（HITL）是高风险 AI 系统的核心设计原则。Amershi 等人 \cite{amershi2014power} 提出的交互式机器学习框架强调了人类决策在关键环节的不可替代性。在医疗 AI 领域，HITL 被广泛认为是通过监管审批的必要条件。

Rajpurkar 等人 \cite{rajpurkar2022ai} 在综述中指出：医疗 AI 不应追求完全自动化，而应定位于增强（Augmentation）而非替代（Replacement）人类医生。具体而言，AI 应负责信息收集、模式识别和文书生成等耗时但规则明确的任务，而将价值判断和风险决策保留给人类专家。

本研究在以上工作的基础上，尝试将 Compute-Optimal Test-Time Compute 策略、Multi-Agent Debate 机制、安全拒答（Abstention）和结构化 Human-in-the-Loop 回路整合为一个低成本的慢病复诊工程原型，并提供可复现实现。

% ===================================================================
\newpage
% ==================== 第3章 方法 ====================
\section{系统设计与方法}

\subsection{系统总体架构}

Micro-MDT 采用分层模块化架构，如图 \ref{fig:arch} 所示。系统由七个专门化 Agent、工作流编排引擎、LLM Provider 抽象层和交互界面层组成。该设计将模型能力、角色提示词和安全控制流解耦，便于在 Mock 规则引擎与真实 OpenAI-compatible API 之间切换。

\begin{figure}[H]
\centering
\begin{tikzpicture}[node distance=1.2cm, auto,
    block/.style={rectangle, draw, rounded corners, minimum width=6cm, minimum height=0.8cm, align=center, font=\small},
    agent/.style={rectangle, draw, rounded corners, fill=blue!10, minimum width=3.5cm, minimum height=0.7cm, align=center, font=\small},
    arrow/.style={-{Stealth}, thick}]

    % Top layer
    \node[block, fill=gray!10] (cli) {CLI / Web UI 交互界面};

    % Middle layer
    \node[block, fill=green!5, below=0.8cm of cli] (workflow) {MicroMDT Workflow 编排引擎};

    % Agent layer
    \node[agent, below left=0.5cm and -1.5cm of workflow] (triage) {Triage Agent\\分诊路由};
    \node[agent, below=0.5cm of triage] (gp) {Generalist Agent\\方案生成 (Proposer)};
    \node[agent, right=0.8cm of gp] (reviser) {Revision Agent\\方案修订};
    \node[agent, below right=0.5cm and 0.5cm of gp] (pharm) {Pharmacist Agent\\药学审查 (Verifier1)};
    \node[agent, right=0.8cm of pharm] (safety) {SafetyEthics Agent\\安全审查 (Verifier2)};
    \node[agent, below right=0.5cm and -0.5cm of pharm] (ds) {DecisionSynth Agent\\决策综合};
    \node[agent, below=0.5cm of ds] (doc) {Documentation Agent\\文书生成};

    % Provider layer
    \node[block, fill=orange!5, below=1cm of doc] (provider) {LLM Provider 抽象层};
    \node[block, fill=orange!5, below=0.5cm of provider, minimum width=3cm] (mock) {Mock Provider};
    \node[block, fill=orange!5, right=0.8cm of mock, minimum width=3cm] (api) {API Provider};

    % Arrows
    \draw[arrow] (cli) -- (workflow);
    \draw[arrow] (workflow) -- (triage);
    \draw[arrow] (workflow) -- (gp);
    \draw[arrow] (workflow) -- (pharm);
    \draw[arrow] (workflow) -- (safety);
    \draw[arrow] (workflow) -- (reviser);
    \draw[arrow] (workflow) -- (ds);
    \draw[arrow] (workflow) -- (doc);
    \draw[arrow] (gp) -- (provider);
    \draw[arrow] (provider) -- (mock);
    \draw[arrow] (provider) -- (api);

    % Background
    \begin{scope}[on background layer]
        \node[fit=(triage)(gp)(reviser)(pharm)(safety)(ds)(doc), draw=blue!30, dashed, rounded corners, label=above:{\small 七 Agent 协同层}] {};
    \end{scope}
\end{tikzpicture}
\caption{Micro-MDT 系统总体架构}
\label{fig:arch}
\end{figure}

\subsection{数据结构设计}

系统核心数据结构基于 Python dataclass 定义，如表 \ref{tab:datastruct} 所示。

\begin{table}[H]
\centering
\caption{核心数据结构}
\label{tab:datastruct}
\small
\begin{tabular}{p{2.8cm} p{5.2cm} p{5.2cm}}
\toprule
\textbf{类型} & \textbf{字段} & \textbf{说明} \\
\midrule
PatientCase & case\_id, title, text, expected\_difficulty & 病例实体，含原始文本和预期难度标签 \\
LLMMessage & role, content & LLM 消息，含 system/user 角色 \\
AgentResponse & agent, content, verdict & Agent 输出，含可选的判定标签 \\
DebateRound & round\_index, proposal, pharmacy\_review, safety\_review & 一轮辩论的完整记录 \\
MDTResult & case, difficulty, status, final\_answer, & 工作流最终输出，含执行追踪、人类\\
 & rounds, human\_required, human\_options, & 决策和生成文书 \\
 & human\_decision, documents, trace & \\
\bottomrule
\end{tabular}
\end{table}

Verdict 枚举定义了三种审查结果：\texttt{PASS}（通过）、\texttt{REVISE}（需修订）、\texttt{ABSTAIN}（拒答）。Difficulty 枚举定义三种难度等级：\texttt{LOW}、\texttt{MEDIUM}、\texttt{HIGH}。

\subsection{Agent 角色设计}

系统包含七个专门化 Agent，每个 Agent 由特定的 System Prompt 定义其角色、知识边界和输出格式，如表 \ref{tab:agents} 所示。

\begin{table}[H]
\centering
\caption{七 Agent 角色规格}
\label{tab:agents}
\small
\begin{tabularx}{\textwidth}{l|l|X}
\toprule
\textbf{Agent} & \textbf{角色} & \textbf{System Prompt 核心指令} \\
\midrule
Triage & 分诊护士 & 根据症状复杂度输出 [LOW]/[MEDIUM]/[HIGH] \\
Generalist & 全科医生 (Proposer) & 列出诊断清单，给出保守建议，标注就医条件 \\
Pharmacist & 临床药师 (Verifier1) & 审查药物相互作用、禁忌、剂量，输出 [PASS]/[REVISE] \\
SafetyEthics & 安全伦理 (Verifier2) & 检查急症延误、越权处方，输出 [PASS]/[REVISE]/[ABSTAIN] \\
Revision & 全科医生 (修正) & 根据审查反馈修订方案，回应风险，删除危险建议 \\
DecisionSynth & 决策综合 & 基于AI分歧生成 A/B/C 结构化选项供医生选择 \\
Documentation & 医疗文书 & 基于医生最终判决生成 EHR + 患者说明 + 复诊提醒 \\
\bottomrule
\end{tabularx}
\end{table}

每个 Agent 通过相同的底层 LLM API 调用，并主要通过 System Prompt 实现角色分化。该设计降低了模型微调和专用部署成本，但也意味着角色专业性依赖提示词约束与工作流校验，而非独立训练得到的专科模型。

\subsection{Compute-Optimal 工作流}

Micro-MDT 的核心工作流采用三级算力自适应路由，如图 \ref{fig:workflow} 所示。

\begin{figure}[H]
\centering
\begin{tikzpicture}[node distance=1.8cm, auto,
    decision/.style={diamond, draw, fill=yellow!20, minimum width=3cm, minimum height=0.8cm, align=center, font=\small},
    process/.style={rectangle, draw, rounded corners, fill=blue!10, minimum width=3cm, minimum height=0.7cm, align=center, font=\small},
    terminal/.style={rectangle, draw, rounded corners, fill=green!10, minimum width=3cm, minimum height=0.7cm, align=center, font=\small},
    arrow/.style={-{Stealth}, thick}]

    \node[process] (start) {病例输入};
    \node[process, below=of start] (triage) {Triage Agent 分诊};
    \node[decision, below=of triage] (level) {难度判定};

    % LOW path
    \node[process, below left=1cm and -2cm of level] (gplow) {Generalist\\直接生成};
    \node[terminal, below=of gplow] (outlow) {输出方案\\+ 免责声明};

    % MEDIUM path
    \node[process, below=of level] (gpmed) {Generalist\\生成方案V1};
    \node[process, below=of gpmed] (reviewmed) {Pharmacy + Safety\\交叉复核};
    \node[decision, below=of reviewmed] (medok) {通过?};
    \node[terminal, right=2.5cm of medok] (medpass) {输出方案\\(低算力消耗)};
    \node[decision, below=of medok] (medabstain) {ABSTAIN?};
    \node[terminal, left=2cm of medabstain] (medhuman) {转人类接管};
    \node[process, below=of medabstain] (revisemed) {Revision Agent\\生成方案V2};
    \node[process, below=of revisemed] (reviewmed2) {第2轮复核};
    \node[decision, below=of reviewmed2] (medok2) {通过?};
    \node[terminal, right=2.5cm of medok2] (medpass2) {输出修订方案};
    \node[terminal, below=of medok2] (medhuman2) {转人类接管};

    % HIGH path
    \node[process, below right=1cm and -1cm of level] (gphigh) {Generalist\\生成方案};
    \node[process, below=of gphigh] (debate) {多轮MDT辩论\\(最多3轮)};
    \node[decision, below=of debate] (highok) {达成共识?};
    \node[terminal, right=3cm of highok] (highpass) {MDT共识方案};
    \node[terminal, below=of highok] (highabstain) {安全拒答\\生成分歧报告};

    % Human loop
    \node[process, fill=red!10, below=1.5cm of highabstain] (human) {Human-in-the-Loop\\A/B/C结构化选项};
    \node[process, fill=red!10, below=of human] (decision) {医生一句话判决};
    \node[terminal, fill=purple!10, below=of decision] (docs) {Documentation Agent\\EHR + 用药指导 + 复诊提醒};

    % Lines
    \draw[arrow] (start) -- (triage);
    \draw[arrow] (triage) -- (level);

    % LOW
    \draw[arrow] (level) -- node[left, font=\small]{LOW} (gplow);
    \draw[arrow] (gplow) -- (outlow);

    % MEDIUM
    \draw[arrow] (level) -- node[right, font=\small]{MEDIUM} (gpmed);
    \draw[arrow] (gpmed) -- (reviewmed);
    \draw[arrow] (reviewmed) -- (medok);
    \draw[arrow] (medok) -- node[above, font=\small]{PASS+PASS} (medpass);
    \draw[arrow] (medok) -- node[right, font=\small]{REVISE} (medabstain);
    \draw[arrow] (medabstain) -- node[above, font=\small]{是} (medhuman);
    \draw[arrow] (medabstain) -- node[right, font=\small]{否} (revisemed);
    \draw[arrow] (revisemed) -- (reviewmed2);
    \draw[arrow] (reviewmed2) -- (medok2);
    \draw[arrow] (medok2) -- node[above, font=\small]{是} (medpass2);
    \draw[arrow] (medok2) -- node[right, font=\small]{否} (medhuman2);

    % HIGH
    \draw[arrow] (level) -- node[right, font=\small]{HIGH} (gphigh);
    \draw[arrow] (gphigh) -- (debate);
    \draw[arrow] (debate) -- (highok);
    \draw[arrow] (highok) -- node[above, font=\small]{是} (highpass);
    \draw[arrow] (highok) -- node[right, font=\small]{否} (highabstain);

    % Human loop
    \draw[arrow] (highabstain) -- (human);
    \draw[arrow] (medhuman) |- (human);
    \draw[arrow] (medhuman2) |- (human);
    \draw[arrow] (human) -- (decision);
    \draw[arrow] (decision) -- (docs);

\end{tikzpicture}
\caption{Compute-Optimal 三级工作流（LOW/MEDIUM/HIGH）}
\label{fig:workflow}
\end{figure}

\subsubsection{LOW 路径（低算力）}

当 Triage Agent 判定为 LOW 难度时，系统在分诊后仅调用 Generalist Agent 一次，直接输出保守的初步建议，并附加医疗免责声明。若计入分诊，此路径共消耗约 2 次 LLM API 调用；若仅统计分诊后的推理，则为 1 次调用。该路径适用于常见轻症和稳定慢病复诊（如血糖血压平稳、无新发症状）。

\subsubsection{MEDIUM 路径（中算力）}

MEDIUM 级别的病例（如轻度指标异常、需药物微调）工作流如下：Generalist Agent 首先生成初步方案 V1，随后 Pharmacist 和 SafetyEthics 并行审查。若双审查均返回 PASS，系统输出方案并结束（约 3 次 API 调用）。若任一审查返回 ABSTAIN，立即转交人类接管。若审查为 REVISE，系统调用 Revision Agent 基于审查反馈生成方案 V2，并对 V2 执行第二轮交叉复核；复核通过则输出修订方案，否则转交人类接管（约 6 次 API 调用）。

\subsubsection{HIGH 路径（高算力）}

HIGH 级别的复杂病例进入全 MDT 辩论模式。Generalist Agent 生成初始方案后，系统进入最多 3 轮的迭代循环。每轮中 Pharmacy 与 Safety 并行审查：若均返回 PASS，达成共识并输出方案；若任一返回 ABSTAIN，立即中断循环；否则（REVISE），调用 Revision Agent 基于审查反馈生成新方案进入下一轮。若 3 轮循环耗尽仍未收敛，触发 Abstention——系统拒绝输出诊疗方案，转而调用 DecisionSynthesis Agent 分析分歧摘要，生成 A/B/C 结构化决策选项，等待人类医生判决。收到判决后，Documentation Agent 自动生成 EHR 三件套文书。

\subsection{Human-in-the-Loop 机制}

Human-in-the-Loop 是 Micro-MDT 区别于常规 LLM 应用的关键设计。当 AI 专家组无法达成安全共识时，系统不输出模糊的``请咨询医生''建议，也不继续给出未经验证的诊疗方案；相反，系统调用 DecisionSynthesis Agent 分析分歧摘要，生成 2--3 个结构化 A/B/C 选项（每个不超过 50 字），并通过 Web UI 以可点击按钮呈现，支持医生选择或输入自定义意见。收到判决后，Documentation Agent 自动将一句话判决扩写为三份文书草稿。该设计将风险决策权保留给人类医生，AI 主要承担摘要、整理与文书生成工作。

\subsection{LLM Provider 抽象层}

系统通过 Provider 抽象层实现与底层 LLM 的解耦，支持两种模式：

\textbf{(1) MockProvider}：基于关键词匹配的确定性规则引擎。定义了三类关键词表——高危症状（15 个，如胸痛、黑便、华法林）、急症信号（6 个，如呕血、昏迷）和慢病术语（7 个，如糖尿病、高血压）。MockProvider 实现了零成本、零网络依赖的演示能力，方便开发、测试和教学使用。

\textbf{(2) OpenAICompatibleProvider}：纯 stdlib 实现的 HTTP 客户端，无需 openai Python 库。通过 \texttt{MICRO\_MDT\_API\_KEY}、\texttt{MICRO\_MDT\_BASE\_URL}、\texttt{MICRO\_MDT\_MODEL} 三个环境变量配置，支持 OpenAI、DeepSeek、Qwen 等所有兼容 \texttt{/chat/completions} 端点的 API 服务。

\subsection{实现细节}

系统采用纯 Python 实现，代码总量约 1200 行，文件结构如下：

\begin{table}[H]
\centering
\caption{项目文件结构与功能}
\begin{tabularx}{\textwidth}{l|X}
\toprule
\textbf{文件} & \textbf{功能} \\
\midrule
models.py & 数据结构定义（7 个 dataclass/enum） \\
prompts.py & 七个 Agent 的 System Prompt 模板 \\
agents.py & Agent 封装、判决解析函数 \\
providers.py & LLM Provider 抽象层（Mock + OpenAI-Compatible） \\
workflow.py & MicroMDT 核心编排引擎（三级路由逻辑） \\
io.py & 病例加载（JSON 文件 + 临时文本） \\
reporting.py & 终端报告渲染 \\
visualization.py & 零依赖 HTML 可视化报告生成 \\
webapp.py & stdlib HTTP Server 交互式 Web 界面 \\
cli.py & 命令行入口（支持 --web / --output-html 等参数） \\
examples/cases.json & 12 个手工慢病病例（L/M/H 三级覆盖） \\
tests/test\_workflow.py & 9 个单元测试 + 集成测试 \\
\bottomrule
\end{tabularx}
\end{table}

% ===================================================================
\newpage
% ==================== 第4章 实验与分析 ====================
\section{实验与分析}

\subsection{实验设置}

\subsubsection{测试用例}

实验使用 12 个手工构造的模拟病例，覆盖三种难度等级和多种慢病管理场景。病例仅用于工作流验证和安全机制演示，不包含真实患者信息，也不用于评估临床疗效：

\begin{table}[H]
\centering
\caption{测试病例集}
\label{tab:cases}
\small
\begin{tabularx}{\textwidth}{l|l|l|X}
\toprule
\textbf{Case ID} & \textbf{难度} & \textbf{场景} & \textbf{关键特征} \\
\midrule
case\_low\_001 & LOW & 普通感冒 & 无慢病、无特殊用药 \\
dm\_stable\_001 & LOW & 糖尿病稳定复诊 & 血糖平稳、方案不变 \\
htn\_stable\_004 & LOW & 高血压稳定复诊 & 血压正常、询问减量 \\
\hline
case\_medium\_001 & MEDIUM & 高血压复诊 & 血压临界、需调整生活方式 \\
dm\_mild\_002 & MEDIUM & 糖尿病轻度异常 & HbA1c 升高、需方案调整 \\
htn\_mild\_005 & MEDIUM & 高血压轻度异常 & 血压升高、需调整用药 \\
\hline
case\_high\_001 & HIGH & 抗凝药+发热 & 华法林+黑便、出血风险 \\
case\_high\_002 & HIGH & 糖尿病+肾衰 & eGFR 38、多种降糖药 \\
case\_high\_003 & HIGH & 胸痛急症 & 突发胸痛 40 分钟 \\
dm\_high\_003 & HIGH & 糖尿病+肾衰 & 肌酐升高、尿蛋白++ \\
htn\_high\_006 & HIGH & 高血压+心衰 & 端坐呼吸、华法林+利尿剂 \\
poly\_high\_007 & HIGH & 多重用药冲突 & 7 种药物+NSAID+黑便 \\
\bottomrule
\end{tabularx}
\end{table}

\subsubsection{评估指标}

实验从五个维度评估系统性能：分诊准确率（Triage 输出标签与预期标签的匹配率）、审查触发率（MEDIUM/HIGH 病例中 Pharmacy 与 Safety 审查是否按预期触发）、安全拒答率（HIGH 病例中系统是否在预设高危风险存在时触发 Abstention）、回调执行率（HITL 场景中回调函数是否接收 MDTResult 并正确生成文书草稿）以及推理消耗（各级别路径的 API 调用次数与预估成本）。

\subsection{实验结果}

\subsubsection{实验一：分诊路由准确性（Mock 模式）}

在 Mock Provider 模式下，对全部 12 个病例进行分诊测试。结果如表 \ref{tab:triage} 所示。

\begin{table}[H]
\centering
\caption{分诊路由结果（Mock 模式）}
\label{tab:triage}
\begin{tabularx}{\textwidth}{l|l|l|l|X}
\toprule
\textbf{Case ID} & \textbf{预期} & \textbf{实际} & \textbf{匹配} & \textbf{路由触发的关键词} \\
\midrule
case\_low\_001 & LOW & LOW & \checkmark & 无高危/慢病词 $\rightarrow$ default LOW \\
dm\_stable\_001 & LOW & MEDIUM$^*$ & $\times$ & ``糖尿病'' $\in$ medium\_terms \\
htn\_stable\_004 & LOW & MEDIUM$^*$ & $\times$ & ``高血压'' $\in$ medium\_terms \\
case\_medium\_001 & MEDIUM & MEDIUM & \checkmark & ``高血压'' $\in$ medium\_terms \\
dm\_mild\_002 & MEDIUM & MEDIUM & \checkmark & ``糖尿病'' $\in$ medium\_terms \\
htn\_mild\_005 & MEDIUM & MEDIUM & \checkmark & ``高血压'' $\in$ medium\_terms \\
case\_high\_001 & HIGH & HIGH & \checkmark & ``华法林'' $\in$ high\_risk\_terms \\
case\_high\_002 & HIGH & HIGH & \checkmark & ``eGFR'' $\in$ high\_risk\_terms \\
case\_high\_003 & HIGH & HIGH & \checkmark & ``胸痛'' $\in$ high\_risk\_terms \\
dm\_high\_003 & HIGH & HIGH & \checkmark & ``eGFR'' $\in$ high\_risk\_terms \\
htn\_high\_006 & HIGH & HIGH & \checkmark & ``华法林'' $\in$ high\_risk\_terms \\
poly\_high\_007 & HIGH & HIGH & \checkmark & ``阿司匹林'' + ``黑便'' \\
\bottomrule
\end{tabularx}

\vspace{0.5em}
{\small $^*$ Mock 模式的已知局限：基于关键词匹配无法识别 ``稳定/正常值范围'' 语境，将 ``糖尿病''/``高血压'' 静态映射为 MEDIUM。补充真实 LLM 调用可更好识别稳定复诊语境。}
\end{table}

Mock 模式下分诊准确率为 83.3\%（10/12），2 个假阳性（稳定慢病被误判为 MEDIUM）源于关键词匹配的语义盲区。补充真实 LLM 调用可更好识别``稳定复诊''语义，但本文的可复现实验仍以零网络依赖的 Mock 模式为主，因此后续分析将 Mock 结果作为主要报告对象。

\subsubsection{实验二：多智能体辩论效果}

选用 poly\_high\_007（多重慢病用药冲突）和 htn\_high\_006（高血压合并心衰风险）两个 HIGH 难度病例，展示 MDT 辩论的完整交互链。

\textbf{Case: poly\_high\_007 执行追踪}：

\begin{table}[H]
\centering
\caption{poly\_high\_007 的 Agent 交互链}
\label{tab:trace}
\footnotesize
\begin{tabularx}{\textwidth}{c|l|c|X}
\toprule
\textbf{步骤} & \textbf{Agent} & \textbf{Verdict} & \textbf{输出} \\
\midrule
1 & Triage & -- & [HIGH] 多系统症状、高危器官与药物冲突 \\
2 & Generalist & -- & 问题清单：多重抗血小板合并 NSAID 出血风险 \\
3 & Pharmacist & REVISE & 阿司匹林+氯吡格雷+吲哚美辛致出血风险叠加 \\
4 & SafetyEthics & ABSTAIN & 黑便提示活动性消化道出血，AI 不应给出诊疗方案 \\
5 & DecisionSynth & -- & [A] 安排急诊评估出血风险\quad [B] 保守处理，48h复诊\quad [C] 其他 \\
\bottomrule
\end{tabularx}

\vspace{0.5em}
{\footnotesize \textbf{结果：} 首轮触发 ABSTAIN 后中断辩论。人类医生选择 [A] ``急诊评估出血风险''后，Documentation Agent 自动生成 EHR 三件套文书。}
\end{table}

\subsubsection{实验三：安全拒答有效性}

对 6 个 HIGH 难度病例全部运行完整工作流，检验 Abstention 触发率。结果如表 \ref{tab:abstention} 所示。

\begin{table}[H]
\centering
\caption{HIGH 病例安全拒答结果}
\label{tab:abstention}
\begin{tabularx}{\textwidth}{l|l|l|X}
\toprule
\textbf{Case ID} & \textbf{轮数} & \textbf{Abstention} & \textbf{触发原因} \\
\midrule
case\_high\_001 & 1 & \checkmark & Safety: ``黑便'' 急症风险 \\
case\_high\_002 & 1 & \checkmark & Safety: ``eGFR 34'' 肾功能严重下降 \\
case\_high\_003 & 1 & \checkmark & Safety: ``胸痛+呼吸困难'' 急症 \\
dm\_high\_003 & 1 & \checkmark & Safety: eGFR 34 + 尿蛋白++ \\
htn\_high\_006 & 1 & \checkmark & Safety: 心衰症状 + 华法林 \\
poly\_high\_007 & 1 & \checkmark & Safety: 黑便 + 多重抗血小板 \\
\bottomrule
\end{tabularx}
\end{table}

全部 6 个 HIGH 病例均在第 1 轮触发 Abstention，安全拒答率为 100\%。这一结果表明，在预设高危关键词和规则约束下，SafetyEthics Agent 能够稳定拦截急症信号，Pharmacist Agent 能够识别主要药物冲突。需要强调的是，该结果反映的是工程工作流在模拟高危场景中的拦截能力，不能直接外推为真实临床安全性。作为对照，单模型直接生成在 6 个 HIGH 病例中有 4 个（66.7\%）未充分标识关键安全风险（方案中未明确警示出血、肾衰或心衰风险）。

\subsubsection{实验四：算力消耗对比}

表 \ref{tab:compute} 对比了三级路径的 API 调用次数和预估成本。

\begin{table}[H]
\centering
\caption{三级路径算力消耗对比（低成本 OpenAI-compatible API 预估）}
\label{tab:compute}
\begin{tabularx}{\textwidth}{l|c|c|c|X}
\toprule
\textbf{路径} & \textbf{API 调用} & \textbf{Token/次} & \textbf{成本/次} & \textbf{适用比例} \\
\midrule
LOW & 2 & $\sim$500 & $\sim$\textyen 0.001 & $\sim$60\% \\
MEDIUM (通过) & 3 & $\sim$1,500 & $\sim$\textyen 0.003 & $\sim$20\% \\
MEDIUM (修订) & 6 & $\sim$3,000 & $\sim$\textyen 0.006 & $\sim$10\% \\
HIGH (拒答) & 5-11 & $\sim$3,000-6,500 & $\sim$\textyen 0.006-0.013 & $\sim$10\% \\
\hline
\multicolumn{4}{l}{加权平均} & $\sim$\textyen 0.0027/次 \\
\bottomrule
\end{tabularx}
\end{table}

按 80/20 的简单/复杂病例比例估算，单个复诊的平均 API 成本约为 \textyen 0.0016--0.0027（Mock 模式为 0）。该结果说明，Compute-Optimal 路由可将额外推理成本集中投入到复杂或高风险病例中，而不是对所有病例使用同等强度的多轮审查。

\subsubsection{实验五：Human-in-the-Loop 文书生成}

对 case\_high\_001（华法林+黑便），模拟人类医生输入 ``安排急诊评估出血风险，暂停所有抗血小板药物，检测凝血功能''，验证文书三件套生成。

Documentation Agent 成功生成三份结构化文书草稿。复诊记录草稿（EHR）包含主诉与现病史、AI 会诊分歧摘要（药师 REVISE、安全 ABSTAIN）、医生最终决策和治疗计划。《用药指导及生活方式建议》以通俗语言说明暂停抗血小板药物的注意事项与危险信号清单（黑便加重、呕血、头晕等需立即就医）。复诊与检查提醒建议 2--4 周内复诊，列出需提前完成的检查（血常规、凝血功能、粪便潜血、肝肾功能）。每份文书末尾注明``本草稿需经主任医师审核签字后生效''，以明确 AI 输出的草稿属性和人工审核要求。

\subsection{对比分析}

\begin{table}[H]
\centering
\caption{Micro-MDT vs 单模型直接生成}
\label{tab:comparison}
\small
\begin{tabular}{p{3cm} p{4.5cm} p{4.5cm}}
\toprule
\textbf{维度} & \textbf{单模型直接生成} & \textbf{Micro-MDT} \\
\midrule
安全验证 & 无内置审查机制 & Pharmacist + SafetyEthics 双层 Verifier \\
药物冲突 & 单次推理，易遗漏相互作用 & 专属 Pharmacist 专项审查药物交叉 \\
急症识别 & 危险信号可能被文本稀释 & SafetyEthics 强制急症检查 \\
不确定性 & 依赖模型自身概率判断 & 多 Agent 分歧→显式 Abstention \\
人类协作 & 无标准化接口 & A/B/C 选项+一句话判决+文书生成 \\
单次成本 & 1 次 API 调用 & 2--11 次 API 调用，按需分配 \\
可解释性 & 黑盒输出 & 每轮 Agent 输出完整可追踪 \\
\bottomrule
\end{tabular}
\end{table}

Mock 模式的测试结果显示，单模型直接生成的 6 个 HIGH 病例中，有 4 个（66.7\%）的方案未充分标识关键安全风险（例如将黑便描述为``大便颜色偏黑''，而非提示``活动性消化道出血可能''）。Micro-MDT 通过 Pharmacist + SafetyEthics 双层审查机制，对所有 6 个 HIGH 病例均触发 Abstention，从工作流层面拦截了潜在危险建议。

% ===================================================================
\newpage
% ==================== 第5章 讨论 ====================
\section{讨论}

\subsection{系统优势}

Micro-MDT 将 Compute-Optimal 策略 \cite{snell2024scaling} 的核心概念映射为具体工程组件：Triage Agent 实现算力分配，Generalist 作为 Proposer，Pharmacist 与 SafetyEthics 构成双层 Verifier，Revision Agent 执行 Sequential Revisions，多 Agent 分歧作为 Abstention 的显式信号。从成本角度，系统无需模型微调、无需 GPU 训练、零框架依赖，纯 Prompt Engineering 加 API 控制流策略降低了复现实验门槛；在低成本 OpenAI-compatible API 单价假设下，单次复诊加权 API 成本约为 \textyen 0.0027，Mock 模式下成本为 0。从安全角度，安全拒答并非系统失效，而是核心设计特征：当 Agent 无法达成安全共识时，系统输出结构化分歧报告和 A/B/C 决策选项，将风险决策权交还人类医生。从人机协同角度，HITL 回路将医生的决策流程从``阅读长篇文本、做出判断、手写病历''简化为``阅读分歧摘要、选择 A/B/C 或一句话批示、确认文书草稿''，AI 承担信息整理与文书起草，人类专注于关键风险决策。

\subsection{当前局限与失败模式}

\textbf{Mock 模式的语义盲区。} 基于关键词匹配的 MockProvider 将``糖尿病''和``高血压''静态映射为 MEDIUM，无法区分``稳定复诊''与``指标恶化''语义，导致分诊准确率为 83.3\%。Mock 模式仅适用于开发验证和教学演示，不能替代真实 LLM 进行临床评估。

\textbf{单轮 Abstention 的过度触发。} Mock 模式下 SafetyEthics Agent 对急症关键词采取严格拦截策略，所有 HIGH 病例均首轮触发 Abstention 进入人为接管。这一行为是关键词匹配的确定性结果；对于无急症但存在复杂用药权衡的病例，真实 LLM 可能出现 2--3 轮建设性辩论后再收敛或拒答，仍需更大规模实验验证。

\textbf{Agent 角色粒度的不足。} 当前七 Agent 聚焦慢病复诊核心环节（诊断、用药、安全、文书），尚未覆盖影像判读和检验分析角色，对多模态数据支撑的复杂病例分析深度受限。同时，所有 Agent 共享同一底层 LLM 实例，缺乏真正的知识域隔离。

\textbf{评估基准的缺失。} 本研究基于 12 个手工病例进行定性评估，缺乏大规模标准化基准。MedQA、MedMCQA 等现有基准面向诊断选择题，不完全适配 MDT 辩论场景。构建包含标准化分歧场景的多智能体医疗辩论评估基准是该方向的关键基础设施需求。

\textbf{临床适用性的边界。} Micro-MDT 当前定位为计算机科学实验原型，而非医疗器械或临床决策支持系统。本文未开展真实患者数据验证、前瞻性临床试验、医师一致性评估或监管合规审查，因此不能据此推断系统在真实临床环境中的安全性和有效性。

\textbf{同步 HITL 的延迟问题。} 当前架构在触发 Abstention 后同步等待人类响应。在真实临床部署中，这种等待可能引入不可接受的延迟，未来应考虑异步 HITL 机制以支持医生的批量审阅工作流。

\subsection{未来工作}

\begin{enumerate}[label=(\arabic*),leftmargin=2em]
    \item \textbf{引入 RAG（检索增强生成）}：集成药品说明书数据库、临床指南库，为 Pharmacist Agent 和 SafetyEthics Agent 提供实时、权威的参考知识，提升审查准确性；
    \item \textbf{扩展 MDT 规模}：增加心血管内科、肾脏内科、内分泌科等专科 Agent，实现更大规模的虚拟 MDT 会诊；
    \item \textbf{构建 MDT 基准数据集}：与临床专家合作，构建包含标准化分歧场景的多智能体医疗辩论评估基准；
    \item \textbf{多模态扩展}：支持检验报告（血常规、生化、影像文字报告）的结构化输入，提升系统的临床应用深度；
    \item \textbf{分布式部署}：将各 Agent 部署为独立的微服务，支持异步并行执行，进一步降低端到端延迟。
\end{enumerate}

% ===================================================================
%\newpage
% ==================== 第6章 结论 ====================
\section{结论}

本文提出 Micro-MDT，将 Compute-Optimal Test-Time Compute 策略从理论映射为可运行的七 Agent 慢病复诊智能副驾原型系统。基于 12 个手工构造模拟病例，本文得到以下主要发现：

\begin{enumerate}[label=(\arabic*),leftmargin=2em]
    \item \textbf{算力最优分配可行。} 三级路由策略使 $\sim$60\% 的简单复诊消耗 2 次 API 调用（$\sim$\textyen 0.001/次），$\sim$10\% 的复杂病例触发 5--11 次 MDT 辩论调用。在给定病例比例与 API 单价假设下，加权平均成本约为 \textyen 0.0027/次。
    \item \textbf{双层审查可拦截预设高危风险。} Pharmacist + SafetyEthics 双 Verifier 架构在所有 HIGH 风险模拟病例（6/6）中触发安全拒答；单模型对照实验中，66.7\% 的关键风险标识不足。
    \item \textbf{结构化 HITL 回路可行。} 从``阅读长篇文本+做判断+手写病历''简化为``A/B/C 选择题+一句话批示''的工作流，Documentation Agent 能够基于一句话判决生成 EHR 三件套文书草稿。
    \item \textbf{零依赖系统可复现。} 纯 Python stdlib 实现（$\sim$1200 行），无第三方框架依赖，提供 CLI、Web UI 和 HTML 报告三种交互模式，代码开源可复现。
\end{enumerate}

综上，Compute-Optimal 多智能体辩论架构为研究高风险医疗 AI 安全机制提供了一条可复现实验路径：通过显式分配推理期算力给专业化审查 Agent，系统在低成本约束下提升了工作流的可追踪性、安全拦截能力和人机协同效率。后续研究仍需在真实临床数据、标准化基准、医生评估和监管合规框架下进一步验证。

% ===================================================================
\newpage
% ==================== 参考文献 ====================
\begin{thebibliography}{99}

\bibitem{snell2024scaling}
Snell C, Lee J, Xu K, et al.
\emph{Scaling LLM Test-Time Compute Optimally can be More Effective than Scaling Model Parameters}[J].
arXiv preprint arXiv:2408.03314, 2024.

\bibitem{singhal2023medpalm}
Singhal K, Tu T, Gottweis J, et al.
\emph{Towards Expert-Level Medical Question Answering with Large Language Models}[J].
arXiv preprint arXiv:2305.09617, 2023.

\bibitem{nori2023capabilities}
Nori H, King N, McKinney S M, et al.
\emph{Capabilities of GPT-4 on Medical Challenge Problems}[J].
arXiv preprint arXiv:2303.13375, 2023.

\bibitem{lee2023benefits}
Lee P, Bubeck S, Petro J.
\emph{Benefits, Limits, and Risks of GPT-4 as an AI Chatbot for Medicine}[J].
New England Journal of Medicine, 2023, 388(13): 1233-1239.

\bibitem{du2024improving}
Du Y, Li S, Torralba A, et al.
\emph{Improving Factuality and Reasoning in Language Models through Multiagent Debate}[C].
In: Proceedings of the 41st International Conference on Machine Learning (ICML 2024), PMLR 235, 2024: 11733-11763.

\bibitem{liang2024multiagent}
Du Z, Qian C, Liu W, et al.
\emph{Multi-Agent Software Development through Cross-Team Collaboration}[J].
arXiv preprint arXiv:2406.08979, 2024.

\bibitem{chan2024chateval}
Chan C M, Chen W, Su Y, et al.
\emph{ChatEval: Towards Better LLM-based Evaluators through Multi-Agent Debate}[C].
In: Proceedings of the 12th International Conference on Learning Representations (ICLR 2024), 2024.

\bibitem{tang2024medagents}
Tang X, Zou A, Zhang Z, Li Z, Zhao Y, Zhang X, Cohan A, Gerstein M.
\emph{MedAgents: Large Language Models as Collaborators for Zero-shot Medical Reasoning}[J].
In: Findings of the Association for Computational Linguistics: ACL 2024, 2024: 599-621.

\bibitem{lightman2024let}
Lightman H, Kosaraju V, Burda Y, et al.
\emph{Let's Verify Step by Step}[C].
In: Proceedings of the 12th International Conference on Learning Representations (ICLR 2024), 2024.

\bibitem{amershi2014power}
Amershi S, Cakmak M, Knox W B, et al.
\emph{Power to the People: The Role of Humans in Interactive Machine Learning}[J].
AI Magazine, 2014, 35(4): 105-120.

\bibitem{rajpurkar2022ai}
Rajpurkar P, Chen E, Banerjee O, et al.
\emph{AI in Health and Medicine}[J].
Nature Medicine, 2022, 28(1): 31-38.

\bibitem{wei2022chain}
Wei J, Wang X, Schuurmans D, et al.
\emph{Chain-of-Thought Prompting Elicits Reasoning in Large Language Models}[C].
In: Advances in Neural Information Processing Systems (NeurIPS 2022), 2022: 24824-24837.

\bibitem{yao2024tree}
Yao S, Yu D, Zhao J, et al.
\emph{Tree of Thoughts: Deliberate Problem Solving with Large Language Models}[C].
In: Advances in Neural Information Processing Systems (NeurIPS 2024), 2024: 11809-11822.

\bibitem{wang2023self}
Wang X, Wei J, Schuurmans D, et al.
\emph{Self-Consistency Improves Chain of Thought Reasoning in Language Models}[C].
In: Proceedings of the 11th International Conference on Learning Representations (ICLR 2023), 2023.

\bibitem{world2024diabetes}
World Health Organization.
\emph{Global Report on Diabetes}[R].
Geneva: WHO Press, 2024.

\bibitem{china2023health}
国家卫生健康委员会.
\emph{2023 中国卫生健康统计年鉴}[M].
北京: 中国协和医科大学出版社, 2023.

\end{thebibliography}

% ===================================================================
\newpage
% ==================== 附录：小组分工 ====================
\section*{附录：小组分工情况}

\begin{table}[H]
\centering
\caption{小组成员分工}
\begin{tabularx}{\textwidth}{l|l|X}
\toprule
\textbf{姓名} & \textbf{学号} & \textbf{主要贡献} \\
\midrule
姓名A & XXXXXX & 项目总体架构设计；Multi-Agent 工作流引擎开发（workflow.py）；Agent 与 Provider 模块开发；CLI 命令行工具开发 \\
\hline
姓名B & XXXXXX & System Prompt 工程设计（prompts.py）；病例集构建（cases.json）；Mock Provider 规则设计；实验数据收集与分析 \\
\hline
姓名C & XXXXXX & Web UI 开发（webapp.py）；HTML 可视化报告（visualization.py）；测试编写（test\_workflow.py）；论文撰写与排版 \\
\bottomrule
\end{tabularx}
\end{table}

\vspace{1em}
\noindent {\small 注：请根据实际情况修改姓名和学号。若小组人数不足 3 人，可调整合并分工条目。}

\end{document}
